% ===== PRÉAMBULE CANONIQUE — à coller VERBATIM en tête de CHAQUE .tex =====
\documentclass[11pt,a4paper]{article}
\usepackage[utf8]{inputenc}\usepackage[T1]{fontenc}\usepackage{lmodern}
\usepackage[french]{babel}
\usepackage{amsmath,amssymb,mathtools}\usepackage{icomma}
\usepackage[version=4]{mhchem}          % \ce{...} : équations & formules
\usepackage{siunitx}\sisetup{locale=FR} % \qty{}{}, \si{}, \ang{}
\usepackage{chemfig}                    % \chemfig{...} : structures développées
\usepackage{xcolor}\usepackage{colortbl}
\usepackage{tikz}\usepackage{pgfplots}\pgfplotsset{compat=1.18}
\usepackage[most]{tcolorbox}\usepackage{enumitem}\usepackage{array,booktabs}
\usepackage[a4paper,margin=2.2cm]{geometry}
\usetikzlibrary{arrows.meta,calc,angles,quotes,positioning,patterns,backgrounds,shapes.geometric}
\definecolor{primary}{RGB}{222,49,99}\definecolor{primarydark}{RGB}{150,22,55}
\definecolor{defcol}{RGB}{0,114,178}\definecolor{methcol}{RGB}{52,92,148}
\definecolor{excol}{RGB}{0,135,90}\definecolor{attncol}{RGB}{200,40,55}
\tcbset{boxbase/.style={enhanced,breakable,boxrule=0.9pt,arc=2.5pt,
  left=3mm,right=3mm,top=2.3mm,bottom=2.3mm,fonttitle=\bfseries,coltitle=white}}
% --- boîtes TUTORIEL (noms EXACTS, ne pas renommer) ---
\newtcolorbox{cle}[1][]{boxbase,colback=defcol!6,colframe=defcol,
  title={\textsc{À retenir}\ifx\relax#1\relax\relax\else\textmd{ : \itshape #1}\fi}}
\newtcolorbox{methode}[1][]{boxbase,colback=methcol!7,colframe=methcol,sharp corners,
  title={\textsc{Méthode}\ifx\relax#1\relax\relax\else\textmd{ --- \itshape #1}\fi}}
\newtcolorbox{exemplebox}[1][]{boxbase,colback=excol!5,colframe=excol,
  title={\textsc{Exemple}\ifx\relax#1\relax\relax\else\textmd{ : \itshape #1}\fi}}
\newtcolorbox{piege}[1][]{boxbase,colback=attncol!6,colframe=attncol,
  title={\textsc{Piège}\ifx\relax#1\relax\relax\else\textmd{ : \itshape #1}\fi}}
\newtcolorbox{remarque}[1][]{boxbase,colback=black!4,colframe=black!45,
  title={\textsc{Remarque}\ifx\relax#1\relax\relax\else\textmd{ : \itshape #1}\fi}}
% --- boîtes EXERCICE / CORRIGÉ (noms EXACTS, auto-comptées) ---
\newtcolorbox[auto counter]{exo}[1][]{boxbase,colback=primary!4,colframe=primary,
  title={\textsc{Exercice}~\thetcbcounter\ifx\relax#1\relax\relax\else\textmd{ --- \itshape #1}\fi}}
\newtcolorbox[auto counter]{corrige}[1][]{boxbase,colback=excol!4,colframe=excol,
  title={\textsc{Corrigé}~\thetcbcounter\ifx\relax#1\relax\relax\else\textmd{ --- \itshape #1}\fi}}
\newcommand{\R}{\mathbb{R}}\newcommand{\N}{\mathbb{N}}\newcommand{\Z}{\mathbb{Z}}\newcommand{\ds}{\displaystyle}
% (un \usepackage{fancyhdr}+en-tête est possible ; il est ignoré côté web)
% =========================================================================
\begin{document}
\begin{exo}[appliquer la formule]
Calcule la charge formelle de l'atome, avec
$\mathrm{C.F.} = N_v - 2\times(\text{doublets libres}) - (\text{liaisons})$.
\begin{enumerate}[label=\alph*)]
  \item \ce{O} : $2$ doublets libres, $2$ liaisons.
  \item \ce{O} : $3$ doublets libres, $1$ liaison.
  \item \ce{N} : $1$ doublet libre, $3$ liaisons.
  \item \ce{N} : $0$ doublet libre, $4$ liaisons.
  \item \ce{C} : $0$ doublet libre, $4$ liaisons.
\end{enumerate}
\end{exo}

\begin{exo}[molécules neutres]
Calcule la charge formelle de l'atome souligné (elles valent toutes $0$ : vérifie-le !).
\begin{enumerate}[label=\alph*)]
  \item \ce{O} dans \ce{H2O} ($2$ doublets libres, $2$ liaisons).
  \item \ce{N} dans \ce{NH3} ($1$ doublet libre, $3$ liaisons).
  \item \ce{C} dans \ce{CH4} ($0$ doublet libre, $4$ liaisons).
  \item \ce{Cl} dans \ce{HCl} ($3$ doublets libres, $1$ liaison).
  \item \ce{H} dans \ce{H2O} ($0$ doublet libre, $1$ liaison).
\end{enumerate}
\end{exo}

\begin{exo}[atomes chargés]
Calcule la charge formelle de l'atome central de ces ions.
\begin{enumerate}[label=\alph*)]
  \item \ce{N} dans \ce{NH4+} ($0$ doublet libre, $4$ liaisons).
  \item \ce{O} dans \ce{H3O+} ($1$ doublet libre, $3$ liaisons).
  \item \ce{O} dans \ce{OH-} ($3$ doublets libres, $1$ liaison).
  \item \ce{B} dans \ce{BF4-} ($0$ doublet libre, $4$ liaisons).
\end{enumerate}
\end{exo}

\begin{exo}[la somme redonne la charge]
Pour chaque édifice, on donne les charges formelles de ses atomes. Vérifie que leur somme est
bien égale à la charge globale.
\begin{enumerate}[label=\alph*)]
  \item \ce{CO3^2-} : $\ce{C}=0$ ; un $\ce{O}=0$ ; deux $\ce{O}=-1$.
  \item \ce{NH4+} : $\ce{N}=+1$ ; quatre $\ce{H}=0$.
  \item \ce{NO3-} : $\ce{N}=+1$ ; un $\ce{O}=0$ ; deux $\ce{O}=-1$.
\end{enumerate}
\end{exo}

\begin{exo}[interpréter le signe]
Un atome porte la charge formelle indiquée. A-t-il, dans cette structure, \emph{plus},
\emph{moins} ou \emph{autant} d'électrons « propres » que l'atome neutre isolé ?
\begin{enumerate}[label=\alph*)]
  \item $\mathrm{C.F.} = +1$.
  \item $\mathrm{C.F.} = -1$.
  \item $\mathrm{C.F.} = 0$.
\end{enumerate}
\end{exo}

\end{document}
